% Generated draft: requires mathematical and source review.
\documentclass[11pt]{article}
\usepackage[margin=25mm]{geometry}
\usepackage{fontspec}
\setmainfont{Latin Modern Roman}
\usepackage{amsmath,amssymb,mathtools,unicode-math}
\setmathfont{Latin Modern Math}
\usepackage{enumitem,needspace,xurl,hyperref}
\setlength{\parindent}{0pt}
\setlength{\parskip}{5pt}
\setlength{\emergencystretch}{3em}
\newcommand{\N}{\mathbb{N}}
\newcommand{\Z}{\mathbb{Z}}
\newcommand{\Q}{\mathbb{Q}}
\newcommand{\R}{\mathbb{R}}
\newcommand{\C}{\mathbb{C}}
\newcommand{\F}{\mathbb{F}}
\newcommand{\A}{\mathbb{A}}
\newcommand{\PP}{\mathbb{P}}
\newcommand{\E}{\mathbb{E}}
\newcommand{\eps}{\varepsilon}
\newcommand{\dd}{\,\mathrm{d}}
\begin{document}
\section*{UnsolvedMath ID: 30002988}
\textbf{Source problem code:} catalogue entry
\subsection{Definitions and mathematical statement}
The supplied question concerns the intermediate regularity range
\[
\frac13<\alpha<\frac23
\]
for the equation
\[
\operatorname{Det}\nabla^2v
=-\frac12\operatorname{curl}\operatorname{curl}(\nabla v\otimes\nabla v)=f.
\]
The source does not specify the domain, the function space indexed by \(\alpha\), the assumptions on \(f\), boundary conditions, or the meaning of the separate term ``weak solution.'' Consequently, the following definitions explain the displayed operator under a planar interpretation; they do not supply missing hypotheses for the original question.

For this interpretation, let \(\Omega\subset\R^2\) be open. For a matrix field \(A=(A_{ij})_{i,j=1}^2\), set
\[
\operatorname{curl}\operatorname{curl}A
:=\partial_{22}A_{11}+\partial_{11}A_{22}
-\partial_{12}A_{12}-\partial_{21}A_{21},
\]
with derivatives understood in the sense of distributions. If \(v\in W^{1,2}_{\mathrm{loc}}(\Omega)\), then \((\nabla v\otimes\nabla v)_{ij}=\partial_i v\,\partial_j v\) is locally integrable, so the displayed very weak Hessian determinant is well defined as a distribution. Writing \(v_i=\partial_i v\), the equation with \(f\in\mathcal D'(\Omega)\) means
\[
-\frac12\int_\Omega
\left(v_1^2\partial_{22}\varphi
+v_2^2\partial_{11}\varphi
-2v_1v_2\partial_{12}\varphi\right)\,\dd x
=\langle f,\varphi\rangle
\quad\text{for every }\varphi\in C_c^\infty(\Omega).
\]
For smooth \(v\), this operator agrees with \(\det\nabla^2v\).

If the intended regularity is \(v\in C^{1,\alpha}_{\mathrm{loc}}(\Omega)\), this means that \(v\) is continuously differentiable and, for every compact \(K\subset\Omega\), there is \(C_K<\infty\) such that
\[
|\nabla v(x)-\nabla v(y)|\le C_K|x-y|^\alpha
\quad\text{for all }x,y\in K.
\]
This interpretation of \(\alpha\) is not explicitly specified in the supplied record and requires confirmation.

\emph{Source question.} Determine the qualitative behavior of weak and very weak solutions of the displayed Monge--Amp\`ere equation in the regime \(1/3<\alpha<2/3\). The source does not identify a particular rigidity, convexity, flexibility, existence, or uniqueness assertion to prove, nor does it specify the quantifiers over domains and right-hand sides. Thus this is a research question with unresolved formulation details, rather than a fully specified conjecture.

The supplied dated assessment describes rigidity and convexity above the two-thirds threshold and flexible constructions below the one-third threshold, while describing the intermediate range as unresolved. These are supplied literature qualifications, not independently verified mathematical assertions; their applicable hypotheses are absent from the record. The endpoints \(\alpha=1/3\) and \(\alpha=2/3\) are excluded from the stated question.
\subsection{Short English statement}
The question asks how weak and very weak Monge--Amp\`ere solutions behave when \(1/3<\alpha<2/3\), between the rigidity and flexibility regimes described by the supplied literature assessment. A precise conjecture still requires the intended regularity space, domain, assumptions on \(f\), solution notions, boundary conditions, and a specific qualitative conclusion. The supplied status is ``partially solved''; this draft does not independently verify that assessment.
\subsection{Sources}
\begin{itemize}
\item Supplied local UnsolvedMath registry record, permanent ID \(30002988\), source problem code \texttt{catalogue entry}, title ``Intermediate Regularity Regime for Monge--Amp\`ere Solutions,'' in the Oberwolfach Reports Open Problems collection. No individual proposer is supplied; the proposed year is recorded as 2015.
\item Supplied report attribution: \emph{Mini-Workshop: Scales in Plasticity (2016)}, Oberwolfach Reports, pp.~3007--3048, DOI \texttt{10.4171/owr/2015/53}. \url{https://doi.org/10.4171/owr/2015/53}. The background describes Oberwolfach Reports as documenting mathematical workshops and their open-problem discussions. The report citation is supplied provenance; its contents and the precise location of this question have not been independently verified.
\item Supplied literature source, ``Recent very weak Monge--Ampere results'': \url{https://link.springer.com/article/10.1007/s00526-025-03019-0}.
\item Supplied literature source, ``Intermediate-regularity analysis'': \url{https://arxiv.org/abs/2310.06693}.
\item Supplied literature source, ``Convex-integration constructions'': \url{https://sites.pitt.edu/~lewicka/preprints/Flexible-Monge9.pdf}.
\item Literature assessment checked at \texttt{2026-08-21}: supplied dated curation, not independently verified. It labels the problem ``partially solved,'' describes rigidity and convexity above the two-thirds threshold and flexible constructions below the one-third threshold, and describes the intermediate range as unresolved. The supplied background calls this a dated literature triage derived from OWR curation and requests independent verification of time-sensitive status and cited claims.
\item The registry supplies the statement-status label \texttt{corrected\_verified}, with the explanation that notation and publication prose were repaired without changing the question's substance. This is a supplied editorial verification claim, not mathematical review of this draft or independent source verification.
\end{itemize}
\end{document}
